\subsubsection{Interpretation of the difference order}
\label{sec:order-interpretation}
For equally spaced samples $x(t_i)$ with spacing $h>0$,
the operator $D_d$ approximates an order-$d$ temporal derivative:
\begin{equation}\label{eq:difference-derivative-approximation}
\frac{\Delta^d x(t)}{h^d}\approx x^{(d)}(t+dh/2).
\end{equation}
The approximation is centered at the middle of the $d$-step
stencil and has second-order accuracy for sufficiently
smooth underlying functions at an interior location.
Thus $d=1$ regularizes changes in level (discrete slope),
$d=2$ regularizes changes in slope (discrete acceleration
for position data), and $d=3$ regularizes
changes in acceleration, or \emph{jerk}.
For example,
\begin{equation}\label{eq:discrete-jerk}
j\!\left(t+\tfrac32h\right)
\approx
\frac{x(t+3h)-3x(t+2h)+3x(t+h)-x(t)}{h^3}.
\end{equation}
If position $x(t)$ is cubic with constant physical jerk $j$,
then $\Delta^3 x=jh^3$ \emph{exactly}. Consequently
$\mu=0$ in the third-difference penalty favors a
quadratic (constant-acceleration) limiting trend, whereas
a prescribed nonzero $\mu=jh^3$ favors a cubic
(constant-jerk) limiting trend. This distinction follows
from the affine constraint $D_3\tau=\mu\mathbf1$,
not from a new physical forecasting mechanism.

The differential interpretation requires correct sampling
units. If $y$ is recorded position, $\Delta^3 y$ is in
position units, and only $\Delta^3 y/h^3$ has jerk units.
Taking finite differences of noisy positions can strongly
amplify measurement noise; the jerk of a fitted trend is an
\emph{estimated} derivative rather than an independent
physical observation \citep{vanbreugel2020}.
The third-order example is therefore illustrative of the
general method, not evidence that $d=3$ is universally
optimal or yields reliable forecasts.

\paragraph{When $d$ is unknown.}
The difference order can itself be treated as a discrete
model-selection hyperparameter, alongside the normalized
smoothness target $S$ and, if appropriate, a prespecified
reference-mean policy. For a finite candidate set
$\mathcal D\subset\{1,\ldots,n-1\}$, a transparent
selection rule is
\begin{equation}\label{eq:joint-order-selection}
(\widehat d,\widehat s)
\in\arg\min_{\substack{d\in\mathcal D\\s\in\mathcal S_d}}
\widehat R_{\rm valid}(d,s),
\end{equation}
where $\widehat R_{\rm valid}$ is a validation loss
computed on data withheld from fitting and tuning.
For future-time prediction, forward-chaining or
rolling-origin validation preserves chronology and
should be distinguished from retrospective in-sample
GCV/LOOCV \citep{tashman2000,bergmeir2018}.
Nested validation or an untouched test period is required
to assess the performance of an adaptively selected
$(d,s)$ without reusing the same validation outcomes
for the final performance claim.

Although $S_d\in[0,1]$ has the same normalized
\emph{trace interpretation} for each $d$,
different orders penalize different derivatives, have
different nullspaces, and generally imply different
selected $\lambda$ and fitted trends. The smallest
validated loss, not the numerical value of $S$
alone, should determine whether a change of order
is justified. This is an extension of the fixed-$d$
analysis in the present theoretical results; no new
multi-order predictive experiment is asserted.
